\section{Method}

A digit-test report begins with a table of nonnegative reported counts and a
chosen unit of analysis: contest, jurisdiction, reporting unit, batch, or a
larger count table. The method extracts a digit feature and compares its
observed distribution to a stated reference distribution. The comparison may be
summarized by a chi-square statistic, an absolute or squared distance, a
ranked-residual score, or another predeclared distance measure. The statistic is
an anomaly score, not a verdict.

\begin{center}
\small
\begin{tabularx}{\textwidth}{lX}
\toprule
Family & Reference and caveat \\
\midrule
First digit & Compare leading digits 1 through 9 with the Benford reference,
$P(d)=\log_{10}(1+1/d)$. Use only where the count table plausibly spans broad
magnitudes; many precinct tables do not. \\
Second digit & Compare digits 0 through 9 with the corresponding Benford-style
second-digit reference. Treat results as especially weak for vote counts,
because modest range restriction and precinct design can move the distribution. \\
Last digit & Compare terminal digits 0 through 9 with a uniform reference, or
with a predeclared administrative baseline if rounding or reporting rules affect
final digits. Last-digit concentration can be a screen, but it can also reflect
aggregation, batch construction, or data-entry conventions. \\
\bottomrule
\end{tabularx}
\end{center}

For an observed count $O_d$ and expected count $E_d$ in digit bin $d$, a common
summary is
\[
  X^2 = \sum_d \frac{(O_d-E_d)^2}{E_d},
\]
provided the expected counts are large enough for the approximation to be
meaningful. A simpler distance score such as $\sum_d |O_d-E_d|$ may be reported
instead when the goal is ranking rather than calibration. If a p-value is
reported, the transcript must state the reference distribution, exclusions, and
baseline population that make the calibration meaningful.

The central diagnostic warning is part of the method. Vote counts often violate
Benford assumptions for innocent reasons: bounded electorates, similar precinct
sizes, contest-specific turnout ceilings, truncation by reporting-unit design,
and pooling choices that mix unlike jurisdictions. A high statistic may indicate
a data feature worth explaining, but it is weak evidence by itself and must be
read with the structure of the count table.
